\section{Methodology}
\label{sec:methodology}

This section describes the algorithms, randomness sources, convergence
mechanisms, and software architecture that together constitute the experimental
framework.  All design decisions are documented in sufficient detail to permit
independent replication.

%% ----------------------------------------------------------------
\subsection{Algorithm Implementations}
\label{sec:meth:algorithms}

Both the classical PSO and the quantum-behaved QPSO are implemented as
self-contained Python classes that share a common interface: each accepts an
objective function, dimensionality, search bounds, and a set of
hyperparameters, and returns a dictionary containing the best position, best
score, convergence histories, and timing information.

\subsubsection{Classical PSO}

The PSO implementation follows the constriction coefficient formulation of
Clerc and Kennedy~\cite{Clerc2002}, with $c_1 = c_2 = 1.494$, which
corresponds to the analytically derived stable regime with
$\varphi = c_1 + c_2 \approx 2.988$ and constriction factor
$\chi \approx 0.7298$.  The inertia weight is decayed linearly from
$w_{\max} = 0.9$ to $w_{\min} = 0.4$ over the iteration budget according to
Equation~\eqref{eq:inertia_decay}.  Velocity clamping is enforced in every
dimension: the maximum permissible speed is set to 20\% of the search range,
i.e.\ $v_{\max} = 0.2\,(x_{\max} - x_{\min})$.  After the velocity update,
each component is clipped to $[-v_{\max},\; v_{\max}]$ before being added to
the position, and the resulting position is itself clipped to the search bounds
to prevent particles from leaving the feasible region.

A \emph{fine-tuning phase} is activated during the final 10\% of the iteration
budget.  In this phase, the inertia weight is overridden to a fixed value of
$0.3$, and the cognitive and social coefficients are both increased to $2.5$.
The rationale is that, by the end of the optimisation, the swarm has
typically converged to a small neighbourhood of the best-known solution, and
the reduced inertia combined with stronger attraction to personal and global
bests encourages intensive local search.  This mechanism is a variant of the
common practice of switching to a local search strategy in the terminal phase
of a metaheuristic run.

\subsubsection{Quantum-Behaved QPSO}

The QPSO implementation follows the canonical formulation of Sun et
al.~\cite{Sun2004,Sun2012} as described in
Section~\ref{sec:bg:qpso}, with the fitness-weighted mean best position of Xi
et al.~\cite{Xi2008} replacing the standard unweighted centroid.  The
contraction--expansion coefficient $\beta$ is decayed linearly from
$\beta_{\max} = 1.0$ to $\beta_{\min} = 0.5$.  When the stagnation detector
indicates that no improvement has occurred for half of the stagnation limit, a
widening factor of 1.3 is applied to $\beta$ (capped at $\beta_{\max}$) to
temporarily broaden the quantum well and facilitate escape from the current
basin.

The position update proceeds in four steps.  First, an attractor point
$\mathbf{p}_i$ is computed as a random interpolation between the particle's
personal best and the global best (Equation~\ref{eq:qpso_attractor}).  Second,
a quantum displacement $\boldsymbol{\delta}_i$ is computed from the distance
between the weighted mean best and the particle's current position, scaled by
$\beta$ and the logarithmic inverse-CDF factor
(Equation~\ref{eq:qpso_displacement}).  Third, the displacement is added to or
subtracted from the attractor according to a binary tunnelling sign.  Fourth,
a wave-packet momentum term is applied: each particle maintains an exponential
moving average of its recent position changes, and a fraction of this momentum
vector is added to the quantum-sampled position.  The momentum update for
particle~$i$ is
\begin{equation}
  \mathbf{m}_i^{t+1} = \gamma_d\,\mathbf{m}_i^{t}
    + (1 - \gamma_d)\,(\mathbf{x}_i^{t+1,\text{quantum}} - \mathbf{x}_i^{t}),
  \label{eq:momentum_update}
\end{equation}
where $\gamma_d = 0.7$ is the momentum decay rate, and the final position is
$\mathbf{x}_i^{t+1} = \mathbf{x}_i^{t+1,\text{quantum}}
+ \gamma_w\,\mathbf{m}_i^{t+1}$ with momentum weight $\gamma_w = 0.3$.  This
mechanism is motivated by the observation that the standard QPSO position
distribution models only the squared modulus $|\psi|^2$ of the wave function,
discarding the momentum information $p = \hbar k$ encoded in the phase.
Restoring a momentum component enables directional persistence that aids
valley-following on functions such as Rosenbrock.

The full QPSO position update is summarised in
Algorithm~\ref{alg:qpso_update}.

\begin{algorithm}[t]
\caption{QPSO position update for particle $i$ at iteration $t$.}
\label{alg:qpso_update}
\begin{algorithmic}[1]
\Require $\mathbf{pbest}_i$, $\mathbf{gbest}$, $\mathbf{mbest}_w$,
         $\mathbf{x}_i^t$, $\mathbf{m}_i^t$, $\beta$
\State $\boldsymbol{\phi} \leftarrow \text{Uniform}(0,1)^D$
\State $\mathbf{p}_i \leftarrow \boldsymbol{\phi} \odot \mathbf{pbest}_i
       + (1 - \boldsymbol{\phi}) \odot \mathbf{gbest}$
       \Comment{Attractor}
\State $\mathbf{u} \leftarrow \text{Uniform}(0.001,\, 0.999)^D$
\State $\boldsymbol{\delta}_i \leftarrow \beta\,
       |\mathbf{mbest}_w - \mathbf{x}_i^t| \odot
       \ln(1/\mathbf{u})$
       \Comment{Quantum displacement}
\State $\mathbf{s} \leftarrow \text{sign drawn from mode}$
       \Comment{$+1$ or $-1$ per dimension}
\State $\mathbf{x}_i^{q} \leftarrow \mathbf{p}_i
       + \mathbf{s} \odot \boldsymbol{\delta}_i$
       \Comment{Quantum step}
\State $\mathbf{m}_i^{t+1} \leftarrow \gamma_d\,\mathbf{m}_i^t
       + (1 - \gamma_d)\,(\mathbf{x}_i^{q} - \mathbf{x}_i^t)$
       \Comment{Momentum EMA}
\State $\mathbf{x}_i^{t+1} \leftarrow \text{clip}\bigl(
       \mathbf{x}_i^{q} + \gamma_w\,\mathbf{m}_i^{t+1},\;
       x_{\min},\; x_{\max}\bigr)$
\end{algorithmic}
\end{algorithm}


%% ----------------------------------------------------------------
\subsection{Quantum Randomness Modes}
\label{sec:meth:modes}

A central contribution of this work is the systematic comparison of QPSO under
three randomness sources, all accessed through a unified interface that
isolates the source of randomness from the optimisation logic.  The three modes
are summarised in Table~\ref{tab:modes}.

\begin{table}[t]
\centering
\caption{Quantum randomness modes and their random value sources.}
\label{tab:modes}
\begin{tabular}{llll}
\toprule
\textbf{Mode} & \textbf{Continuous values}
  & \textbf{Tunnelling sign}
  & \textbf{Backend} \\
\midrule
\texttt{math}   & NumPy PRNG   & NumPy PRNG
  & None (pure classical) \\
\texttt{hybrid} & NumPy PRNG   & Qiskit $R_Y$ circuit
  & AerSimulator \\
\texttt{full}   & Qiskit $H$ circuits & Qiskit $R_Y$ circuit
  & AerSimulator \\
\bottomrule
\end{tabular}
\end{table}

In \textbf{\texttt{math} mode}, every random quantity---the attractor
interpolation weights $\boldsymbol{\phi}$, the displacement magnitude inputs
$\mathbf{u}$, and the tunnelling sign---is generated by NumPy's Mersenne
Twister pseudo-random number generator.  This mode serves as the baseline
against which the quantum modes are compared.  It requires no quantum software
dependencies and executes fastest among the three modes.

In \textbf{\texttt{full} mode}, all random values are generated by quantum
circuits executed on the Qiskit AerSimulator.  For each particle--dimension
pair, three circuits are constructed per iteration.  The first two are
single-qubit Hadamard circuits: a qubit initialised in the $|0\rangle$ state
is passed through a Hadamard gate $H$ and measured, producing outcome $|1\rangle$
with probability exactly $1/2$.  By repeating this circuit for a fixed number
of \emph{shots} (measurement repetitions) and computing the fraction of
$|1\rangle$ outcomes, a quantum-random floating-point value in $(0,1)$ is
obtained with granularity $1/\text{shots}$.  These two values serve as the
cognitive and social random weights $r_1$ and $r_2$.  The third circuit applies
a Hadamard gate followed by a parameterised $R_Y(\beta\pi)$ rotation:
\begin{equation}
  R_Y(\theta) = \begin{pmatrix}
    \cos(\theta/2) & -\sin(\theta/2) \\
    \sin(\theta/2) & \phantom{-}\cos(\theta/2)
  \end{pmatrix}.
  \label{eq:ry_gate}
\end{equation}
After the combined $H \cdot R_Y(\beta\pi)$ operation, the probability of
measuring $|1\rangle$ is
\begin{equation}
  P(|1\rangle) = \sin^2\!\biggl(\frac{\beta\pi}{2}\biggr).
  \label{eq:prob_one}
\end{equation}
When $\beta$ is large (early exploration phase), $P(|1\rangle) \approx 0.5$,
yielding an approximately unbiased coin flip.  As $\beta$ decays toward 0.5,
$P(|1\rangle) \approx 0.15$, biasing the sign toward $+1$ (outward
displacement) and thereby tightening the effective search radius.  The measured
fraction of $|1\rangle$ outcomes simultaneously provides the displacement
magnitude $u$ and, via thresholding at 0.5, the binary tunnelling sign.  All
circuits for a single iteration are batched into a single job submitted to the
backend, minimising communication overhead.

In \textbf{\texttt{hybrid} mode}, only the tunnelling sign is drawn from a
quantum circuit---the same $H + R_Y(\beta\pi)$ circuit described above---while
all continuous random variables ($\boldsymbol{\phi}$, $\mathbf{u}$, $r_1$,
$r_2$) are generated classically by NumPy.  The motivation for this split is
that shot-based quantum measurements produce values with a granularity of
$1/\text{shots}$.  With a typical shot count of 256, this yields a step size of
approximately 0.004---too coarse for continuous quantities that enter
multiplications and logarithms but entirely appropriate for a binary decision.
The hybrid mode therefore restricts quantum circuits to the decision for which
they are best suited: a genuinely binary random variable with a
$\beta$-dependent bias.


%% ----------------------------------------------------------------
\subsection{Convergence Mechanisms}
\label{sec:meth:convergence}

Both PSO and QPSO incorporate several mechanisms designed to prevent premature
convergence and to accelerate exploitation once the swarm has localised the
global basin.

\subsubsection{Mean Best Position}

The mean best position $\mathbf{mbest}$ serves as the consensus reference for
the swarm.  In classical PSO, it is computed as the unweighted arithmetic mean
of all personal bests (Equation~\ref{eq:mbest}).  In QPSO, the
fitness-weighted variant of Xi et al.\ (Equation~\ref{eq:wqpso_mbest}) is used
instead.  The mean best position enters the QPSO displacement calculation
directly: the quantum step size for each particle is proportional to its
distance from $\mathbf{mbest}$, so particles far from the consensus naturally
take larger steps (exploration) while particles near the consensus take smaller
steps (exploitation).  This self-regulating behaviour is a distinctive
advantage of the QPSO formulation over classical PSO, where exploration versus
exploitation is controlled solely through the global inertia parameter.

\subsubsection{Stagnation Detection and Worst-Particle Scatter}

A stagnation detector monitors the global best score across iterations.  If no
improvement exceeding a numerical tolerance of $10^{-10}$ is observed for a
configurable number of consecutive iterations (the stagnation limit, set to 15
by default), the detector triggers a recovery action: the worst-performing 30\%
of particles---ranked by their personal best fitness---are re-initialised to
random positions within a 15\% neighbourhood of the current global best.
Their personal best scores are reset to infinity, forcing them to re-evaluate
the landscape from fresh vantage points without disturbing the best solution
found so far.  This mechanism provides a convergence guarantee in the weak
sense: the swarm can never remain permanently trapped in a suboptimal region,
because the periodic injection of new search points ensures ongoing
exploration.

\subsubsection{Elitist Local Refinement (QPSO)}

Every five iterations, the QPSO algorithm applies a greedy local refinement to
the global best position using two complementary strategies.  The first
strategy is a \emph{differential mutation step}: two particles are selected at
random, and the normalised difference between their personal bests defines a
search direction.  The global best is perturbed along this direction by a step
size proportional to $\beta$ and the search range.  On valley-type landscapes,
particles tend to spread along the ridge, so the differential vector naturally
aligns with the valley curvature, enabling the refinement to follow the
valley floor.  The second strategy is a \emph{coordinate-wise Gaussian
perturbation}: each dimension of the global best is independently perturbed by
a zero-mean Gaussian with standard deviation $\sigma = 0.02\,\beta\,(x_{\max}
- x_{\min})$, and the perturbation is accepted if and only if it improves the
objective value.  Both strategies are greedy---improvements are incorporated
immediately---and together they provide a local search capability that
compensates for the inherently stochastic, non-directional nature of the
quantum position update.

\subsubsection{Fine-Tuning Phase (PSO)}

As noted in Section~\ref{sec:meth:algorithms}, the classical PSO switches to a
fine-tuning regime in the final 10\% of the iteration budget.  The inertia
weight is reduced to 0.3 and the cognitive and social coefficients are both
increased to 2.5, producing strong attraction toward the personal and global
bests with minimal velocity carry-over.  This configuration aggressively
exploits the neighbourhood of the current best solution and is analogous to the
elitist refinement in QPSO, albeit implemented through the velocity update
mechanism rather than through direct perturbation.


%% ----------------------------------------------------------------
\subsection{Software Architecture}
\label{sec:meth:software}

The framework is implemented as a Python package (\texttt{qpso\_research})
with a modular architecture that separates concerns into distinct modules.
The particle classes (\texttt{PSOParticle}, \texttt{QPSOParticle}) encapsulate
position, velocity (or momentum), personal best tracking, and fitness
evaluation.  The convergence utilities---mean best computation, weighted mean
best computation, stagnation detection, and worst-particle scatter---reside in
a shared module used by both optimisers.  The quantum engine is implemented as
a standalone class that exposes two methods: \texttt{sample\_batch\_full()} for
\texttt{full} mode and \texttt{sample\_signs\_only()} for \texttt{hybrid}
mode.  The benchmark functions and their metadata (bounds, known optima) are
registered in a dictionary that maps string names to dataclass instances,
enabling configuration-driven experiment specification.

A key design decision is the \emph{lazy import} of Qiskit.  The quantum engine
module is imported only when the QPSO class is instantiated with
\texttt{mode="full"} or \texttt{mode="hybrid"}.  In \texttt{math} mode, no
quantum software is loaded, allowing the framework to run on systems where
Qiskit is not installed.  This pattern reduces startup time and dependency
footprint for users who wish to use only the classical algorithms.

Experiments are configured through YAML files that specify the benchmark
functions, dimensionalities, number of independent runs, swarm size, iteration
budget, and all algorithm-specific hyperparameters.  A single configuration
file defines an entire experimental campaign, and the experiment runner
iterates over all combinations of benchmark, dimensionality, and algorithm
mode.  Reproducibility is ensured through explicit seed control: the NumPy
random seed is set at the beginning of each run, and the seed value is recorded
alongside the results.  While the Qiskit AerSimulator is itself deterministic
given a fixed seed, the interaction between NumPy and Qiskit random states is
managed carefully to ensure that repeated executions of the same configuration
file produce identical results.

The results of each run---convergence histories, final scores, timing
information, and iteration-to-success counts---are collected into structured
data objects that support both tabular export (for statistical analysis) and
programmatic access (for visualisation).  Statistical comparison utilities
provide non-parametric hypothesis tests and effect-size calculations, enabling
rigorous assessment of whether observed differences between algorithm--mode
combinations are statistically significant.
